\textbf{Het experiment:}

Ons experiment is nagenoeg volledig geautomatiseerd. Bij het runnen van het Python script \linkscript{py} \textit{(Link werkt niet met de gevraagde submap)} zullen er 4 devices aangemaakt worden, elk verbonden met dezelfde router, alles met een max bandwidth van 200Kbps.\\

\textbf{Het script runnen:}
Run het script door: "python3 L2-10-1.py" uit te voeren\\
In mininet run: "py run\_test(net)"
Wat er achterliggend gebeurt:
- h1 is de server, en runt 3 iperf3 servers.\\
- h2 en h3 zijn tcp clients die elk verbinden met een server van h1\\
- h4 is een udp client die verbindt met een server van h1\\
De outputfiles wordt automatisch opgeslagen in een folder op de desktop van de VM.

\textbf{Aanwezige files in de traces:}
\begin{itemize}
    \item Main wiresharkbestand "capture.pcapng" (uit het perspectief van h1)
    \item Bestanden met de client outputs:
    \begin{itemize}
        \item h2-tcp.txt
        \item h3-tcp.txt
        \item h4-udp.txt
    \end{itemize}
\end{itemize}

\textbf{Conclusies:}
-h1 en h2 beginnen alleen, ze verzadigen de verbinding snel.\\
Nadat h3 erbij komt wordt TCP-fairness toegepast om de gelimiteerde bandbreedte van h1 eerlijk te verdelen tussen -h2 en h3.\\
-Zodra h4 erbij komt, neemt deze zonder aandacht voor de andere verbindingen alle bandbreedte in, TCP wordt weggeconcureerd door de 'onvriendelijkere' UDP.\\